% Distinct supplemental primary: Gotama, Nyaya-sutra, Book I,
% Vidyabhusana's English translation/commentary (1913), printed pp.1--21.
% NOT ED26-SRC-063-0105: Ganeri's complete 2004 chapter remains open.
% The corrupt Sanskrit OCR and wider Buddhist/Jaina corpus are not credited.
\chapter{Nyaya Inquiry, Inference, and Debate in the 1913 Translation}
\label{sec:p13-nyaya-sutra-inference-and-debate}
\chaptermark{Nyaya Inquiry and Debate}
\index{Nyaya}\index{Gotama}\index{Vidyabhusana, Satisa Chandra}
\index{inference!Nyaya}\index{debate!Nyaya}

Gotama's \emph{Nyaya-sutra} connects inquiry, knowledge, argument and debate
within a philosophical account of release. This presentation follows Book I,
chapters I--II, in Satisa Chandra Vidyabhusana's English translation and
explanations, printed by the Panini Office in 1913
\cite[Book I, chapters I--II, pp.~1--21]{GotamaVidyabhusana1913NyayaSutras}.
It is one identified Nyaya presentation, not a definition of every Indian
logical tradition. Its admitted testimony, metaphysical background and
argument examples must not be silently attributed to Buddhist or Jaina authors.

The edition prints translated aphorisms followed by explanations, examples,
alternative translations and references to commentators, sometimes within
square brackets. These layers have different attribution. An aphorism's
definition, an English commentator's illustration, and a present account of
the illustration's inferential structure are distinguished below. Locators
such as 1.1.32 mean Book I, chapter I, aphorism 32; 1.2.9 means Book I,
chapter II, aphorism 9. The English numbering is not uniformly continuous:
chapter II initially restarts it and later also prints continuous numbers
beside chapter-local ones. The book/chapter/aphorism locator identifies the
passage. The source of this chapter is the English translation and explanations,
not a newly established Sanskrit critical text.

\section{The categories and the aim of inquiry}
\label{sec:p13-nyaya-categories}

The first aphorism says that knowledge of the true nature of sixteen
categories leads to supreme felicity. The explanation distinguishes their
\emph{enunciation}, \emph{definition} and \emph{critical examination}:
Book I supplies the first two, while the later books examine them. Its aim
is wider than a list of truth-preserving formal rules. The categories are:

\begin{description}[style=nextline,leftmargin=0pt]
 \item[Means of right knowledge (pramana).] The ways by which knowledge
 is claimed to be obtained.
 \item[Object of right knowledge (prameya).] What the inquiry concerns.
 \item[Doubt (samsaya).] An unresolved conflicting judgment about a
 thing's precise character.
 \item[Purpose (prayojana).] What is to be attained or avoided through action.
 \item[Familiar instance (drstanta).] An example agreed upon by an ordinary
 person and an expert.
 \item[Established tenet (siddhanta).] A commitment with the specified
 school, hypothetical or implied standing.
 \item[Members (avayava).] The parts of the stated argument.
 \item[Confutation (tarka).] Reasoning through the alleged absurdity of
 contrary possibilities.
 \item[Ascertainment (nirnaya).] Determination through the reasons on
 opposing sides and removal of doubt.
 \item[Discussion (vada).] Truth-directed disputation.
 \item[Wrangling (jalpa).] Victory-directed defence and attack.
 \item[Cavil (vitanda).] Attack without establishing one's own position.
 \item[Fallacy of reason (hetvabhasa).] A reason failing in a separately
 defined way.
 \item[Quibble (chala).] Opposition through substituted meaning.
 \item[Futility (jati).] An objection resting on mere similarity or
 dissimilarity without the needed connection.
 \item[Occasion for rebuke (nigrahasthana).] The specified failure to
 understand or respond properly in disputation.
\end{description}

The arrangement distinguishes a cognition's source, its object, the uncertainty
motivating inquiry, the articulated argument and the conduct of the exchange.
A correctly arranged argument therefore does not by itself establish that its
reason is known, its example warrants the connection, or its exchange seeks truth.

Aphorism 1.1.2 gives the succession misapprehension, faults, activity, birth and
pain, and says that destruction in the reverse dependence leads to release.
The explanation identifies faults as affection, aversion and stupidity.
Knowledge removes misapprehension; with it the faults, activity and transmigration
cease, so that pain ceases. It also distinguishes what is to be abandoned, what
produces it, its destruction, and the means of that destruction. These are this
account's philosophical commitments, not an empirical study of a logic course.

\subsection{Objects of knowledge and philosophical background}

Aphorism 1.1.9 names soul, body, senses, sense-objects, intellect, mind, activity,
faults, transmigration, fruit, pain and release as objects of right knowledge.
The explanation also mentions the Vaishesika categories substance, quality,
action, generality, particularity, intimate relation and non-existence. That
comparison does not identify the classifications. The ensuing definitions
retain the Nyaya account's own commitments
\cite[1.1.9--22, pp.~5--7]{GotamaVidyabhusana1913NyayaSutras}.

Desire, aversion, volition, pleasure, pain and intelligence are marks of soul
(1.1.10); the explanation calls them qualities of the substance soul.
Body is the site of gesture, senses and sentiments (1.1.11): it strives
toward the desired and away from the hated; senses function in dependence
on its condition; pleasure and pain are experienced there.
Nose, tongue, eye, skin and ear are the five element-produced senses (1.1.12).
Earth, water, light, air and ether are the elements (1.1.13).
The explanation associates nose with earth, tongue with water, eye with light,
skin with air and ear with ether. Smell, taste, colour, touch and sound are
their respective objects and elemental qualities (1.1.14).
These are attributed historical classifications, not contemporary sensory physics.

The text identifies intellect, apprehension and knowledge rather than making
them three different things (1.1.15). The comparison says that Sankhya separates
intellect, evolved from primordial matter, from knowledge in soul; it also notes
that apprehension is commonly used for perception.
The mark of mind is that more than one act of knowledge does not arise at a
time (1.1.16). The explanation requires a conjunction of mind as well as
sense-object contact. With its atomic-mind assumption it argues that mind cannot
be conjoined with more than one sense-organ at once. The conclusion belongs to
those assumptions; it is not a modern experiment on simultaneous processing.

Activity sets voice, mind and body into action (1.1.17). The explanation divides
bodily, vocal and mental actions into good and bad. Its bad bodily cases are
killing, stealing and adultery; its good ones giving, protecting and serving.
Its bad vocal cases are lying, harsh speech, slander and frivolous talk;
its good ones truth, useful and pleasant speech, and reading sacred books.
Its bad mental cases are malice, covetousness and scepticism; its good ones
compassion, freedom from covetousness and devotion. These are the edition's
evaluative examples, not a claim that every act called scepticism is a logical error.

Faults cause activity (1.1.18). Transmigration means repeated births and deaths
(1.1.19): birth connects soul with body, senses, mind, intellect and sentiments,
while death separates them. Fruit is produced by activity and faults (1.1.20)
and is explained as pleasure or pain. Pain is uneasiness (1.1.21), with an
alternative translation describing hindrance to the soul. Release is absolute
deliverance from pain (1.1.22). The explanation rejects identifying release
with an additional eternal pleasure: pleasure, pain and body are impermanent.
The liberation argument retains these distinctions instead of becoming a
numerical utility-maximization claim.

\section{Perception, inference, comparison and testimony}
\label{sec:p13-nyaya-knowledge-sources}

Aphorism 1.1.3 enumerates perception, inference, comparison and word, or verbal
testimony. Bracketed commentary compares other schools' admitted means:
perception alone for Carvakas; perception and inference for Vaishesikas and
Buddhists; these plus testimony for Sankhyas; further presumption for Prabhakaras;
non-existence as a sixth for Bhattas and Vedantins; and probability and rumour
as seventh and eighth for Pauranikas. These are the translator's comparative
claims, not the aphorism's enumeration or full presentations of those schools
\cite[1.1.3--8, pp.~2--4]{GotamaVidyabhusana1913NyayaSutras}.

\subsection{Perception and an alternative rendering}

The principal rendering of 1.1.4 defines perception as knowledge from
sense-object contact, called determinate, unnameable and non-erratic.
The distant observer unable to decide between smoke and dust lacks the
determinate cognition intended. Unnameable is explained as independence
from the thing's name. Apparent water in a heated mirage supplies the
non-erratic qualification: an appearance does not thereby establish water.

The edition prints an alternative rendering in brackets. There non-erratic
cognition may be indeterminate, illustrated by an unspecific recognition
of something, or determinate, illustrated by recognition of a Brahmana.
The two renderings must not be merged into a definition that both requires
determination and permits its absence in the same sense.

\subsection{Three forms of inference}

Aphorism 1.1.5 calls inference knowledge preceded by perception and distinguishes
three kinds. The translator's labels are \emph{a priori}, \emph{a posteriori}
and \emph{commonly seen}. Here a priori does not mean independence from
experience: the explanation begins from a perceived cause.

\begin{description}[style=nextline,leftmargin=0pt]
 \item[From perceived cause to effect.] Seeing clouds and inferring rain.
 \item[From perceived effect to cause.] Seeing a swollen river and inferring
 preceding rain.
 \item[From commonly observed association.] The bracketed cases infer a tail
 from horns, and fire from smoke on a hill.
\end{description}

The explanation reports Vatsyayana's alternative interpretation of the last
kind, as inference to something not commonly seen: affection, aversion and
other qualities are marks from which soul as their substance is inferred.
The alternatives remain distinct, rather than all examples being one causal
rule. The definition gives no statistical error bound for clouds, rivers or
animals, and does not authorize the converse of every sign relation.

\subsection{Comparison and a name's significance}

Comparison (1.1.6) knows something through similarity to something previously
known. A forester says that a \emph{bos gavaeus} is like a cow. Seeing a
cow-like animal in the forest and recalling the instruction, the person
recognizes the animal under that name. The output concerns name and animal
class, not merely perceived shape.
The explanation records the objection that noticing likeness is perception
and needs no separate means. Its reply distinguishes perceiving likeness
from knowing what the name signifies. Whether the reply succeeds turns on
that distinction; the seen animal alone does not identify the proposed outputs.

\subsection{The condition on verbal testimony}

Word (1.1.7) is the instructive assertion of a reliable person. The explanation
uses expertise and willingness to communicate experience. It expressly includes
a mleccha as well as a rishi or Arya, rather than restricting reliability to
one group. The river case has an experienced local elder without enmity toward
the questioner answer whether the river is fordable. Acceptance uses the stated
competence and communication, not a bare assertion's occurrence.

Aphorism 1.1.8 distinguishes testimony about seen and unseen matters. The
explanation says the former is actually verifiable and the latter can be
ascertained by inference where the hearer cannot verify directly.
Its examples are a physician's butter/strength statement and a religious
teacher's horse-sacrifice/heaven statement. These are historical classification
examples, not nutritional advice or newly proved religious reward.
The authority of scripture remains disputed in the debate case below.

\section{Doubt, purpose, familiar instances and tenets}
\label{sec:p13-nyaya-inquiry-conditions}

Doubt (1.1.23) is conflicting judgment about a thing's precise character.
The explanation distinguishes five sources
\cite[1.1.23--31, pp.~7--10]{GotamaVidyabhusana1913NyayaSutras}.

\begin{enumerate}
 \item A common property leaves alternatives unresolved: tallness at
 twilight belongs both to a person and to a post.
 \item A property outside the comparison classes leaves classification
 open: the commentary asks whether sound is eternal when soundness
 belongs neither to its non-eternal animal examples nor eternal atoms.
 \item Conflicting testimony: one philosophical school affirms soul,
 while another denies it.
 \item Irregular perception: actual tank-water and mirage-water appearance
 raise the question whether appearance requires existence.
 \item Irregular non-perception: unperceived water in a radish and absent
 water on dry ground raise the corresponding question about absence.
\end{enumerate}

Purpose (1.1.24) is what one acts to attain or avoid; collecting fuel to cook
is the bracketed example. A familiar instance (1.1.25) is agreed upon by an
ordinary person and an expert. A kitchen with smoke and fire is the example
used for the hill. Agreement on that case, the connection proposed through it,
and its application to the hill remain separately articulated.

An established tenet rests on a school's authority, hypothesis or implication
(1.1.26). Aphorisms 1.1.27--31 distinguish four kinds. A tenet of all schools
is claimed by at least one and opposed by none; the commentary lists the
five elements and sense-objects. A school-specific tenet is accepted by
similar schools but rejected by opposing ones: the Sankhya refusal of
production from nothing and Mimamsa eternity of sound are its examples.
These attributed examples do not establish one shared sound theory.

A hypothetical tenet brings further commitments when accepted. Soul's
recognition of one object by sight and touch, despite being distinct from
the senses, presupposes multiple senses and their particular objects,
knowledge through them, and a substance distinct from and serving as
abode of qualities. An implied tenet emerges from examination without
being explicitly declared: sound's being substance is presupposed in the
eternal/non-eternal debate, and mind's being a sense-organ follows from
other particulars in the explanation. Neither category permits an arbitrary
unstated premise inserted to guarantee the desired conclusion.

\section{Five-member argument and conditions on a reason}
\label{sec:p13-nyaya-five-members}

Aphorism 1.1.32 gives proposition, reason, example, application and conclusion.
The subsequent definitions give their distinct functions
\cite[1.1.32--39, pp.~10--13]{GotamaVidyabhusana1913NyayaSutras}.

\begin{description}[style=nextline,leftmargin=0pt]
 \item[Proposition.] Declaration of what is to be established (1.1.33),
 exemplified by sound's non-eternity.
 \item[Reason.] Establishment through the affirmative or negative
 character of the example (1.1.34--35).
 \item[Example.] A familiar instance has or lacks the property in question
 and exhibits the supplied invariable connection (1.1.36--37).
 \item[Application.] The thesis-subject is said to be so or not so with
 reference to the example (1.1.38).
 \item[Conclusion.] Restatement after the reason; the explanation also
 specifies the example (1.1.39).
\end{description}

\subsection{Hill, smoke and kitchen}

The complete bracketed case at 1.1.32 is:

\begin{enumerate}
 \item The hill has fire: proposition.
 \item The hill has smoke: reason.
 \item Whatever has smoke has fire, as in a kitchen: example.
 \item This hill is so, namely smoky: application.
 \item The hill therefore has fire: conclusion.
\end{enumerate}

The commentary says some add five more matters. One asks whether fire is
claimed throughout the hill or in a part. Another asks whether the supposed
smoke is vapour. A third questions the example's warrant, contrasting a
smoky kitchen with a red-hot iron ball having fire without smoke.
Another identifies the purpose: whether the hill is to be approached,
avoided or regarded with indifference. Finally doubts concerning the hill's
smoke and the invariable connection are to be dispelled. These are reported
additional proposals, not a silent enlargement of the aphorism's five to ten.

Smoke is the sign; fire is to be established. The iron ball shows fire
without smoke, not smoke without fire. It therefore must not be taken
as a counterexample to the differently directed smoke-to-fire connection.
Agreement on the kitchen also does not decide whether the hill's appearance
is smoke. This account of the recorded challenges does not attribute a
modern Boolean calculus to the ancient text.

\subsection{Affirmative sound and pot}

The affirmative case at 1.1.34 and 1.1.36 has all five members:

\begin{enumerate}
 \item Sound is non-eternal.
 \item Sound is produced: the reason.
 \item What is produced is non-eternal, as a pot.
 \item Sound is of that character: it is produced.
 \item Sound is therefore non-eternal.
\end{enumerate}

The pot instantiates the specified produced/non-eternal connection; its
mere resemblance to sound in some other respect would not perform that task.
The thesis depends on the supplied connection and the subject's possession
of the reason. Arrangement supplies no evidence independent of those commitments.

\subsection{Negative sound and soul}

At 1.1.35 and 1.1.37 the negative example says what is not non-eternal is
not produced, as soul. The explanation treats production and eternity as
incompatible. Under those commitments the case is:

\begin{enumerate}
 \item Sound is not eternal.
 \item Sound is produced.
 \item Whatever is eternal is not produced, as soul.
 \item Sound is not of the unproduced character described in the example:
 the given reason identifies it as produced.
 \item Sound is consequently not eternal.
\end{enumerate}

The fourth and fifth lines are editorially corrected formulations, not
unchanged quotations. At 1.1.38 the print's negative application says
"sound is not produced", contrary to its preceding reason; at 1.1.39 its
worked conclusion says "sound is produced", rather than restating the
non-eternity proposition. Here application agrees with reason and conclusion
restates proposition, as the definitions require
\cite[pp.~12--13]{GotamaVidyabhusana1913NyayaSutras}.
This exposes the inference under its supplied premises; it does not prove
those premises for another school, alter the preserved scan or emend Sanskrit.

\section{Confutation and ascertainment}
\label{sec:p13-nyaya-tarka-ascertainment}

Confutation (1.1.40) examines an unknown character by showing the alleged
absurdity of its contrary. The note at 1.1.1 lists confutation, argumentation,
reductio ad absurdum and hypothetical reasoning as English equivalents of
\emph{tarka}; the explanatory use cannot be fixed by one familiar label alone.

Its case asks whether soul is eternal. The explanation says a non-eternal
soul could not enjoy its actions' fruits, transmigrate or attain release,
while those are accepted possibilities. Eternity is thus accepted.
The argument uses those rebirth/action/release commitments, not a purely
logical contradiction without philosophical premises
\cite[1.1.40--41, p.~13]{GotamaVidyabhusana1913NyayaSutras}.

Ascertainment (1.1.41) determines the question and removes doubt through
hearing opposing sides. Opposition alone is not enough: reasons vindicate
a side. The explanation denies that doubt precedes every ascertainment;
perception, accepted scripture or discussion may ascertain directly.
In investigation doubt precedes the process. That inquiry sequence is not
a definition of every possible cognition.

\section{Discussion, wrangling and cavil}
\label{sec:p13-nyaya-debate-forms}

Discussion (1.2.1) adopts an opposing side, articulates it in five members,
defends it by admitted knowledge means, assails its contrary by confutation,
and preserves established tenets. Bracketed commentary distinguishes
truth-directed discussion, victory-directed wrangling and fault-finding cavil
\cite[1.2.1--3, pp.~14--15]{GotamaVidyabhusana1913NyayaSutras}.

The printed discussion case opposes soul's existence and non-existence.
Its affirmative side states existence, gives being an abode of consciousness
as reason, says the non-existent is not an abode of consciousness, exemplified
by a hare's horn, applies the contrary character to soul as abode of
consciousness, and concludes existence. The opponent's corresponding reason
is left incomplete. It must not be invented as an allegedly supplied
historical opponent's proof.

The affirmative participant also invokes scripture, recognition of one
object through sight and touch, and coherence with accepted eternal things
and fruits of actions. The commentary warns that discussion with a Buddhist
opponent rejecting scriptural authority and eternal things would be longer.
The warning matters: consequence under proposed tenets differs from showing
that an opponent shares them.

Wrangling (1.2.2) employs the foregoing defence and attack together with
quibbles, futilities and rebuke-worthy processes for victory. The explanation
describes a speaker concerned with claiming an active successful part,
regardless of whether the arguments support the speaker's own contention.
Cavil (1.2.3) attacks the other side without establishing its own thesis.
These classify purpose and conduct; they do not make every criticism fallacious.

\section{Five failures of a reason}
\label{sec:p13-nyaya-fallacious-reasons}

Aphorism 1.2.4 names erratic, contradictory, equal to the question, unproved
and mistimed reasons. Their following cases expose different mechanisms
\cite[1.2.4--9, pp.~15--17]{GotamaVidyabhusana1913NyayaSutras}.

\subsection{Erratic reason and pervasion}

An erratic reason leads to more than one conclusion (1.2.5).
The commentary displays two opposed attempted five-member cases.
The first is:
\begin{enumerate}
 \item Sound is eternal.
 \item Sound is intangible.
 \item Whatever is intangible is eternal, as atoms.
 \item Sound is so, namely intangible.
 \item Therefore sound is eternal.
\end{enumerate}
The second is:
\begin{enumerate}
 \item Sound is non-eternal.
 \item Sound is intangible.
 \item Whatever is intangible is non-eternal, as intellect.
 \item Sound is so, namely intangible.
 \item Therefore sound is non-eternal.
\end{enumerate}
The universal claims in the third lines are the attempted connections
being challenged, not conclusions endorsed by their appearance here.
The explanation identifies no universal connection between the middle and
major term: the middle is not pervaded by the major. In its example,
intangibility is universally connected neither with eternity nor with
non-eternity. One chosen resemblance does not establish the invariable
connection. Possession of a sign and its adequate connection to the
property to be established are distinct conditions.

\subsection{Contradictory reason}

A contradictory reason opposes the thesis (1.2.6): a pot is claimed
produced because it is eternal, while the example's background excludes
production of the eternal. Accepting that reason supports the contrary
character, unlike a reason merely lacking a known connection.

\subsection{Reason equal to the question}

A reason equal to the question revives the issue it was meant to settle
(1.2.7). Sound's non-eternity is allegedly established by its lacking
eternity. The commentary identifies the content as the same disputed
content: new wording supplies no independent basis.

\subsection{Unproved reason}

An unproved reason needs proof as the proposition does (1.2.8).
Shadow's being a substance is supported by alleged motion. The
commentary points out that motion may instead belong to the person
obstructing light. Until shadow's motion is established, the reason
does not settle its substancehood. That failure does not itself prove
a contrary thesis about every shadow.

\subsection{Mistimed reason and the alternative reading}

The principal explanation of a mistimed reason (1.2.9) concerns when
the proposed connection holds. A jar's colour is manifest through union
with a lamp, and a drum's sound through union with a rod. The attempted
analogy infers sound's durability from colour's existence before and
after illumination. The explanation locates colour's manifestation
during union but sound's at a subsequent moment when union has ceased,
so it rejects the analogy. This historical temporal premise is not a
new acoustical measurement.

The edition also records a reading concerning wrong order of the members,
such as reason before proposition. It reports Vatsyayana's objection:
separation in a Sanskrit sentence can retain legitimate meaning-connection,
and proximity cannot create missing connection. Misordering instead
belongs to an occasion for rebuke called inopportune presentation.
The two interpretations and classifications remain separate.

\section{Quibble and the intended question}
\label{sec:p13-nyaya-quibble}

Quibble opposes through another assumed meaning (1.2.10).
Three kinds concern a term, a genus and a metaphor (1.2.11)
\cite[1.2.10--17, pp.~17--19]{GotamaVidyabhusana1913NyayaSutras}.

The term case (1.2.12) uses \emph{nava-kambala}: the intended boy with
a new blanket is deliberately recast as a boy with nine blankets,
then challenged because he has only one. Ambiguity alone is not the
failure; substitution for the intended claim is.

The genus case (1.2.13) substitutes a broad-class claim for an individual
claim. Learning and conduct are attributed to one Brahmana. An objector
reads that as inferred for everyone in the class and invokes children
lacking them. The commentary's original claim concerns the particular
person with opportunity for study, not a guarantee from class membership.
Its scope contrast is not a present social-class generalization.

The metaphor case (1.2.14) exchanges literal and secondary meaning.
The scaffolds crying out refers to people upon them; the objector
knowingly reads inanimate scaffolds and denies their ability to cry.
The change of reference, not a contradiction in the intended claim,
is the recorded defect.

Aphorisms 1.2.15--17 retain an objection and answer. An objector would
merge metaphor and term quibble because both substitute meaning.
The answer distinguishes direct/literal meaning from the direct/secondary
contrast, and says that similarity alone cannot erase difference:
the same move would collapse all three kinds into one.
This is the source's classification argument, not a general theorem
that every two similar error categories must differ.

\section{Futility, rebuke and the Book I boundary}
\label{sec:p13-nyaya-futility-rebuke}

Futility (1.2.18) rests on mere similarity or dissimilarity.
The explanation requires the pertinent universal connection
\cite[1.2.18--20, pp.~20--21]{GotamaVidyabhusana1913NyayaSutras}.

In the first case, soul is called inactive as all-pervading, like ether.
The opponent invokes resemblance to a pot as seat of union and asks
why soul is not active. The explanation accepts all-pervasion/inactivity
within this account, but no corresponding seat-of-union/activity connection.
The substituted resemblance lacks what carried the original inference.
These metaphysical connections are supplied premises, not proof by
resemblance alone.

The second case invokes dissimilarity: sound is non-eternal because it is
a product unlike ether. The opponent invokes its being an auditory object
unlike a pot to argue for eternity. The explanation links being a product
with exclusion of eternity, but finds no corresponding link between
auditory perception and exclusion of non-eternity. Another dissimilarity
cannot take over the task merely by being a dissimilarity.

An occasion for rebuke arises through misunderstanding or failure to
understand (1.2.19). One example classifies what is not quality as substance,
yet offers a reason admitting only qualities such as colour; the reason
opposes the position. Another has a speaker deny fire's heat, face contrary
evidence of touch, and then deny before the audience having made the claim.
The explanation considers further counterargument unavailing against
persistent misunderstanding or such renunciation. Its harsh descriptions
are historical commentary; identifying the failure does not require using
that language toward a present interlocutor.

Aphorism 1.2.20 says these kinds are multiple. The explanation announces
twenty-four futilities and twenty-two occasions for rebuke for Book V.
Book I supplies the foregoing definitions and cases, not that later full
subclassification. Nor is its five-member argument a modern axiomatic
calculus with an asserted adequacy metatheorem. The link to knowledge
retains the reason, accepted connection, example, subject's character
and philosophical commitments actually supplied.
